% Wersja źródłowa: audyt II, jednostka 15

\section{Pochodne wyższych rzędów i lokalne twierdzenia analizy nieliniowej}
\label{sec:wyzsze-pochodne15}

\providecommand{\atlasref}[2]{#2}


\subsection{Odwzorowania wieloliniowe i definicja wyższych pochodnych}

\noindent\emph{Powiązanie tematyczne w Bibliotece matematycznej:} \atlasref{LA-08::def:lin8-wieloliniowe}{LA-08: odwzorowania wieloliniowe}.


W podstawowym rachunku różniczkowym różniczkę $Df(x)$ traktowaliśmy jako ograniczony operator liniowy. Przy dalszym różniczkowaniu otrzymujemy odwzorowania wieloliniowe. Ta perspektywa pozwala sformułować twierdzenie Taylora i twierdzenia o funkcji odwrotnej oraz uwikłanej bez redukowania przestrzeni Banacha do współrzędnych.

\begin{definicja}[Ciągłe odwzorowania wieloliniowe]
\label{def:multilinear15}
Niech $m\ge1$, $\K\in\{\R,\C\}$, a $E_1,\ldots,E_m,F$ będą przestrzeniami unormowanymi nad $\K$. Przez $\mathcal L^m(E_1,\ldots,E_m;F)$ oznaczamy przestrzeń odwzorowań $A:E_1\times\cdots\times E_m\to F$, liniowych w każdym argumencie i ciągłych, wyposażoną w normę
$\|A\|=\sup\{\|A(h_1,\ldots,h_m)\|:\|h_j\|\le1,\ 1\le j\le m\}$. Jeśli wszystkie dziedziny są równe $E$, piszemy $\mathcal L^m(E;F)$, a dla $m=1$ używamy dotychczasowego oznaczenia $\mathcal L(E,F)$. Przez $\mathcal L^m_s(E;F)$ oznaczamy podprzestrzeń odwzorowań symetrycznych, tj. niezmienniczych przy permutacjach argumentów.
\end{definicja}

\begin{stwierdzenie}[Kryterium ograniczoności i zupełność]
\label{prop:multi-banach15}
Niech $m\ge1$, $\K\in\{\R,\C\}$, $E_1,\ldots,E_m$ będą przestrzeniami unormowanymi nad $\K$, a $F$ przestrzenią Banacha nad tym samym ciałem. Odwzorowanie $m$-liniowe $A:E_1\times\cdots\times E_m\to F$ jest ciągłe wtedy i tylko wtedy, gdy istnieje $C\ge0$ takie, że
$\|A(h_1,\ldots,h_m)\|\le C\prod_{j=1}^m\|h_j\|$ dla wszystkich $h_j\in E_j$. Przestrzeń $\mathcal L^m(E_1,\ldots,E_m;F)$ jest przestrzenią Banacha.
\end{stwierdzenie}
\begin{proof}
Jeśli $A$ jest ciągłe w zerze, istnieje $\delta>0$ takie, że $\|A(v_1,\ldots,v_m)\|\le1$ dla $\|v_j\|<\delta$. Dla niezerowych $h_j$ podstawiamy $v_j=\delta h_j/(2\|h_j\|)$, co po wykorzystaniu wieloliniowości daje oszacowanie z $C=(2/\delta)^m$. Gdy któryś argument jest zerowy, nierówność jest oczywista. Odwrotnie, ustalmy $(h_1,\ldots,h_m)$. Po rozwinięciu wieloliniowym różnica $A(h_1+k_1,\ldots,h_m+k_m)-A(h_1,\ldots,h_m)$ jest sumą wyrazów indeksowanych przez niepuste podzbiory $J\subseteq\{1,\ldots,m\}$. Norma wyrazu o indeksie $J$ nie przekracza $C\prod_{j\in J}\|k_j\|\prod_{j\notin J}\|h_j\|$. Każdy z tych skończenie wielu wyrazów dąży do zera, gdy $\max_j\|k_j\|\to0$. Otrzymujemy ciągłość w każdym punkcie.

Niech $(A_n)$ będzie ciągiem Cauchy'ego w normie operatorowej. Dla ustalonych $h_1,\ldots,h_m$ ciąg $A_n(h_1,\ldots,h_m)$ jest Cauchy'ego w $F$, więc definiuje $A(h_1,\ldots,h_m)$. Granica zachowuje wieloliniowość. Ponieważ normy $\|A_n\|$ są ograniczone przez wspólną stałą, $A$ spełnia oszacowanie wieloliniowe. Przechodząc do granicy w nierówności $\|(A_n-A_p)(h_1,\ldots,h_m)\|\le\varepsilon\prod\|h_j\|$ dla $n,p\ge N$, otrzymujemy $\|A_n-A\|\le\varepsilon$. To dowodzi zupełności.
\end{proof}

\begin{lemat}[Utożsamienie operatorów wieloliniowych]
\label{lem:currying15}
Niech $m\ge1$, $\K\in\{\R,\C\}$, a $E,F$ będą przestrzeniami unormowanymi nad $\K$. Naturalne odwzorowanie
$\mathcal L(E,\mathcal L^m(E;F))\to\mathcal L^{m+1}(E;F)$, zadane przez $T\mapsto\widetilde T$ oraz $\widetilde T(h_0,h_1,\ldots,h_m)=(Th_0)(h_1,\ldots,h_m)$, jest izometrycznym izomorfizmem liniowym.
\end{lemat}
\begin{proof}
Liniowość, wieloliniowość i ciągłość wynikają z oszacowania $\|\widetilde T(h_0,\ldots,h_m)\|\le\|T\|\prod_{j=0}^m\|h_j\|$. Ponadto, z definicji norm operatorowych,
\[
\begin{aligned}
\|\widetilde T\|&=\sup_{\|h_0\|\le1}\sup_{\substack{\|h_j\|\le1\\1\le j\le m}}\|(Th_0)(h_1,\ldots,h_m)\|\\
&=\|T\|.
\end{aligned}
\]
Dla $A\in\mathcal L^{m+1}(E;F)$ definiujemy $(T_Ah_0)(h_1,\ldots,h_m)=A(h_0,h_1,\ldots,h_m)$. Wieloliniowość $A$ daje liniowość $T_A$, zaś $\|T_Ah_0\|\le\|A\|\|h_0\|$ zapewnia jego ograniczoność. Ponieważ $\widetilde T_A=A$, otrzymujemy surjektywność, a izometryczność daje także iniektywność.
\end{proof}

\begin{definicja}[Klasa $C^m$ i pochodne wyższych rzędów]
\label{def:ck15}
Niech $m\ge1$, $\K\in\{\R,\C\}$, $E,F$ będą przestrzeniami Banacha nad $\K$, a $U\subseteq E$ zbiorem otwartym. Wszystkie pochodne Frécheta rozpatrujemy nad $\K$ (dla $\K=\C$ są zespolenie liniowe). Mówimy, że $f:U\to F$ jest klasy $C^m$, gdy $f\in C^1$ oraz rekurencyjnie $D^{j-1}f:U\to\mathcal L^{j-1}(E;F)$ jest klasy $C^1$ dla $2\le j\le m$. Pochodną $D^jf(x)$ utożsamiamy przez Lemat~\ref{lem:currying15} z elementem $\mathcal L^j(E;F)$. Przyjmujemy $D^0f=f$ oraz $C^\infty(U,F)=\bigcap_{m\ge1}C^m(U,F)$. Dla $1\le j\le m$ ciągłość $D^jf$ rozumiemy w normie $\mathcal L^j(E;F)$.
\end{definicja}

\begin{przyklad}[Odwzorowanie wieloliniowe i wielomiany]
\label{ex:multi-poly15}
Niech $E,F$ będą rzeczywistymi przestrzeniami Banacha, $m\ge1$, $B\in\mathcal L^m(E;F)$ i $P:E\to F$ będzie dane wzorem $P(x)=B(x,\ldots,x)$. Wówczas $P\in C^\infty(E,F)$ oraz
$DP(x)h=\sum_{j=1}^m B(x,\ldots,x,\underset{j}{h},x,\ldots,x)$.
Istotnie, różnica $P(x+h)-P(x)$ składa się z wyrazów, w których $h$ występuje co najmniej raz; wyrazy z dokładnie jednym $h$ tworzą wskazany operator, a suma pozostałych jest $O(\|h\|^2)$ przy ustalonym $x$. Indukcyjnie pochodna rzędu $j\le m$ jest sumą po wszystkich iniekcjach $\{1,\ldots,j\}\to\{1,\ldots,m\}$: w pozycjach wskazanych przez iniekcję umieszczamy kolejno $h_1,\ldots,h_j$, a w pozostałych $x$. Różniczkując w następnym kierunku, zastępujemy jeden z pozostałych czynników $x$ przez $h_{j+1}$. Wyrazy z co najmniej dwoma nowymi przyrostami dają resztę $O(\|h_{j+1}\|^2)$ jednostajnie dla $\|h_i\|\le1$, z oszacowania wieloliniowego. To uzasadnia indukcję w normie operatorowej. Ciągłość każdej pochodnej wynika z tego samego oszacowania. Pochodna rzędu $m$ jest stała, a pochodne od rzędu $m+1$ są zerowe. Dla symetrycznego $B$ zachodzi $DP(x)h=mB(x,\ldots,x,h)$, a $D^mP(x)=m!B$.
\end{przyklad}

\subsection{Symetria drugiej pochodnej i reguły rachunku}

\begin{twierdzenie}[Symetria drugiej pochodnej]
\label{thm:schwarz15}
Niech $E,F$ będą rzeczywistymi przestrzeniami Banacha, $U\subseteq E$ otwarty i $f\in C^2(U,F)$. Dla każdego $x\in U$ i $h,k\in E$ zachodzi
$D^2f(x)[h,k]=D^2f(x)[k,h]$. W szczególności $D^2f(x)\in\mathcal L_s^2(E;F)$.
\end{twierdzenie}
\begin{proof}
Dla małych rzeczywistych $s,t$ połóżmy
$\Delta(s,t)=f(x+sh+tk)-f(x+sh)-f(x+tk)+f(x)$. Z dwukrotnego zastosowania wzoru całkowego na odcinku dla funkcji klasy $C^1$ otrzymujemy
\[
\Delta(s,t)=st\int_0^1\int_0^1 D^2f(x+\alpha sh+\beta tk)[k,h] \,\mathrm d\beta\,\mathrm d\alpha.
\]
Istotnie, najpierw piszemy $\Delta(s,t)=s\int_0^1(Df(x+\alpha sh+tk)-Df(x+\alpha sh))h\,\mathrm d\alpha$, a następnie stosujemy wzór na przyrost do $Df$ w kierunku $tk$ i ewaluujemy na $h$. Zgodnie z utożsamieniem z Lematu~\ref{lem:currying15} pierwszy argument $D^2f$ jest kierunkiem różniczkowania $Df$. Jeśli zaczniemy od kierunku $k$, otrzymamy natomiast $st\int_0^1\int_0^1D^2f(x+\alpha sh+\beta tk)[h,k]\,\mathrm d\alpha\,\mathrm d\beta$. Nie trzeba zakładać symetrii ani zamieniać kolejności całkowania, aby porównać granice obu wyrażeń. Dla $s,t\ne0$ dzielimy przez $st$ i przechodzimy z $(s,t)$ do $(0,0)$. Ciągłość $D^2f$ w normie operatorowej zapewnia jednostajną zbieżność funkcji podcałkowych na kwadracie $[0,1]^2$ odpowiednio do $D^2f(x)[k,h]$ i $D^2f(x)[h,k]$. Obie granice są równe, co kończy dowód.
\end{proof}


\begin{wniosek}[Symetria wszystkich pochodnych wyższych rzędów]
\label{cor:symmetric-higher15}
Niech $E,F$ będą rzeczywistymi przestrzeniami Banacha, $U\subseteq E$ otwarty, $m\ge2$ i $f\in C^m(U,F)$. Wówczas $D^mf(x)\in\mathcal L_s^m(E;F)$ dla każdego $x\in U$.
\end{wniosek}
\begin{proof}
Stosujemy indukcję po $m$. Dla $m=2$ jest to Twierdzenie~\ref{thm:schwarz15}. Jeśli $m\ge3$, to odwzorowanie $D^{m-2}f:U\to\mathcal L^{m-2}(E;F)$ jest klasy $C^2$, a więc zastosowanie tego twierdzenia do $D^{m-2}f$ daje symetrię $D^mf(x)$ względem pierwszych dwóch argumentów. Z założenia indukcyjnego $D^{m-1}f(z)$ jest symetryczne dla każdego $z\in U$. Różniczkując tożsamości wyrażające tę symetrię w kierunku pierwszego argumentu, stwierdzamy, że $D^mf(x)$ jest symetryczne także względem ostatnich $m-1$ argumentów. Łącznie otrzymujemy niezmienniczość przy wszystkich zamianach sąsiednich argumentów, a te generują wszystkie permutacje.
\end{proof}

\begin{stwierdzenie}[Druga pochodna złożenia]
\label{prop:second-chain15}
Niech $E,F,G$ będą rzeczywistymi przestrzeniami Banacha, $U\subseteq E$ i $V\subseteq F$ otwarte, $f\in C^2(U,F)$, $g\in C^2(V,G)$ oraz $f(U)\subseteq V$. Wówczas $g\circ f\in C^2(U,G)$ i dla $x\in U$ oraz $h,k\in E$ zachodzi
\[
\begin{aligned}
D^2(g\circ f)(x)[h,k]
={}&D^2g(f(x))[Df(x)h,Df(x)k]\\
 &+Dg(f(x))\bigl(D^2f(x)[h,k]\bigr).
\end{aligned}
\]
\end{stwierdzenie}
\begin{proof}
Z reguły łańcuchowej $D(g\circ f)(x)=Dg(f(x))\circ Df(x)$. Kompozycja operatorów jest ciągłym odwzorowaniem dwuliniowym $\mathcal L(F,G)\times\mathcal L(E,F)\to\mathcal L(E,G)$. Stosując regułę różniczkowania ciągłego odwzorowania dwuliniowego do odwzorowań $x\mapsto Dg(f(x))$ i $x\mapsto Df(x)$, otrzymujemy po zadziałaniu na $k$ dwa składniki. Pierwszy, dzięki regule łańcuchowej zastosowanej do $Dg$, ma postać $D^2g(f(x))[Df(x)h,Df(x)k]$, a drugi jest równy $Dg(f(x))(D^2f(x)[h,k])$. Ciągłość wszystkich występujących pochodnych dowodzi przynależności złożenia do $C^2$.
\end{proof}

\begin{przyklad}[Hesjan funkcjonału na przestrzeni Hilberta]
\label{ex:hessian15}
Niech $H$ będzie rzeczywistą przestrzenią Hilberta, $A\in\mathcal L(H)$ i $b\in H$. Dla $J(x)=\frac12\langle Ax,x\rangle-\langle x,b\rangle$ mamy
$DJ(x)h=\langle h,(A+A^*)x/2-b\rangle$ oraz
$D^2J(x)[h,k]=\langle h,(A+A^*)k/2\rangle$.
Przy konwencji liniowości iloczynu skalarnego w pierwszym argumencie wzór jest symetryczny dla rzeczywistego $H$, ponieważ $(A+A^*)/2$ jest samosprzężony. Jeśli $A=A^*$, to $D^2J(x)[h,k]=\langle h,Ak\rangle$.
\end{przyklad}

\subsection{Twierdzenie Taylora w przestrzeniach Banacha}

\noindent\emph{Powiązanie tematyczne w Bibliotece matematycznej:} \atlasref{AM1-23::ma1:thm:taylor-peano}{AM1-23: twierdzenie Taylora z resztą Peana}.


\begin{twierdzenie}[Wzór Taylora z resztą całkową]
\label{thm:taylor15}
Niech $E,F$ będą rzeczywistymi przestrzeniami Banacha, $U\subseteq E$ otwarty, $m\ge1$, $f\in C^m(U,F)$, a $x\in U$ i $h\in E$ takie, że odcinek $\{x+th:0\le t\le1\}$ zawiera się w $U$. Wtedy
\[
 f(x+h)=\sum_{j=0}^{m-1}\frac1{j!}D^jf(x)[h^j]
 +\frac1{(m-1)!}\int_0^1(1-t)^{m-1}D^mf(x+th)[h^m] \,\mathrm dt,
\]
gdzie $D^0f(x)[h^0]=f(x)$, a $D^jf(x)[h^j]$ oznacza $D^jf(x)[h,\ldots,h]$. Ponadto
$f(x+h)=\sum_{j=0}^{m}\frac1{j!}D^jf(x)[h^j]+r_m(h)$, gdzie
\[
\|r_m(h)\|\le \frac{\|h\|^m}{m!}\sup_{0\le t\le1}\|D^mf(x+th)-D^mf(x)\|.
\]
W szczególności $\|r_m(h)\|=o(\|h\|^m)$ dla $h\to0$.
\end{twierdzenie}
\begin{proof}
Ustalmy $x,h$ i zdefiniujmy $\gamma(t)=f(x+th)$. Z reguły łańcuchowej i indukcji mamy $\gamma^{(j)}(t)=D^jf(x+th)[h^j]$ dla $1\le j\le m$. Pierwsza pochodna jest ciągła, zatem stosujemy podstawowe twierdzenie rachunku całkowego dla funkcji o wartościach w $F$. Udowodnimy wzór dla krzywej $\gamma$ przez indukcję po $m$. Dla $m=1$ jest to twierdzenie Newtona--Leibniza. W kroku od $m$ do $m+1$ zakładamy $\gamma\in C^{m+1}$ i w reszcie rzędu $m$ całkujemy przez części: dla każdego $\ell\in F'$ stosujemy skalarny wzór do $\ell(\gamma^{(m)}(t))$ i funkcji $-(1-t)^m/m!$, której pochodna wynosi $(1-t)^{m-1}/(m-1)!$. Otrzymujemy
\[
 \frac1{(m-1)!}\int_0^1(1-t)^{m-1}\gamma^{(m)}(t)\,\mathrm dt
 =\frac{\gamma^{(m)}(0)}{m!}
 +\frac1{m!}\int_0^1(1-t)^m\gamma^{(m+1)}(t)\,\mathrm dt .
\]
Równość zachodzi w $F$, ponieważ ciągłe funkcjonały przechodzą przez całkę Riemanna i rozdzielają punkty. Podstawienie do wzoru indukcyjnego daje przypadek $m+1$. Wreszcie wstawiamy wyrażenia na $\gamma^{(j)}$.

Odejmujemy od wyrażenia pod całką stałą $D^mf(x)[h^m]$, której całka z wagą $(1-t)^{m-1}/(m-1)!$ jest równa $D^mf(x)[h^m]/m!$. Następnie stosujemy normę operatorową, oszacowanie normy całki przez całkę normy oraz $\int_0^1(1-t)^{m-1}\,\mathrm dt=1/m$. Ciągłość $D^mf$ w $x$ zapewnia, że supremum w oszacowaniu reszty dąży do zera przy $h\to0$.
\end{proof}

\begin{wniosek}[Wzór drugiego rzędu]
\label{cor:taylor2-15}
Niech $E,F$ będą rzeczywistymi przestrzeniami Banacha, $U\subseteq E$ zbiorem otwartym, $f\in C^2(U,F)$ oraz $x\in U$, $h\in E$ takie, że $[x,x+h]\subseteq U$. Wówczas zachodzi
\[
 f(x+h)=f(x)+Df(x)h+\tfrac12D^2f(x)[h,h]
 +\int_0^1(1-t)\bigl(D^2f(x+th)-D^2f(x)\bigr)[h,h] \,\mathrm dt.
\]
Jeżeli $M\ge0$ i dla każdego $z\in[x,x+h]$ zachodzi $\|D^2f(z)\|\le M$, to reszta po wyrazach rzędu pierwszego ma normę co najwyżej $M\|h\|^2/2$.
\end{wniosek}

\begin{stwierdzenie}[Warunki konieczne i wystarczające minimum]
\label{prop:extremum15}
Niech $E$ będzie rzeczywistą przestrzenią Banacha, $U\subseteq E$ otwarty, $x_0\in U$, a $f\in C^2(U,\R)$. Jeżeli $x_0$ jest lokalnym minimum $f$, to $Df(x_0)=0$ i $D^2f(x_0)[h,h]\ge0$ dla każdego $h\in E$. Odwrotnie, jeżeli $Df(x_0)=0$ oraz istnieje $c>0$ takie, że $D^2f(x_0)[h,h]\ge c\|h\|^2$ dla wszystkich $h\in E$, to $x_0$ jest ścisłym lokalnym minimum $f$; dokładniej, dla dostatecznie małych $h$ zachodzi $f(x_0+h)\ge f(x_0)+(c/4)\|h\|^2$.
\end{stwierdzenie}
\begin{proof}
Dla dowolnego $h$ rozważamy rzeczywistą funkcję $\varphi(t)=f(x_0+th)$. Ma ona minimum w $0$, więc $\varphi'(0)=Df(x_0)h=0$ i $\varphi''(0)=D^2f(x_0)[h,h]\ge0$. W drugim kierunku ciągłość $D^2f$ daje promień $r>0$ taki, że $\|D^2f(x_0+u)-D^2f(x_0)\|\le c/2$ dla $\|u\|<r$. Stąd $D^2f(x_0+u)[h,h]\ge(c/2)\|h\|^2$. Z całkowego wzoru Taylora drugiego rzędu, przy $Df(x_0)=0$, otrzymujemy $f(x_0+h)-f(x_0)\ge(c/2)\|h\|^2\int_0^1(1-t)\,\mathrm dt=c\|h\|^2/4$ dla $\|h\|<r$.
\end{proof}

\begin{przyklad}[Dodatniość bez jednostajności nie wystarcza]
\label{ex:positive-hessian15}
W rzeczywistej przestrzeni Hilberta $H=\ell^2(\N;\R)$ niech
$Q(x)=\frac12\sum_{n=1}^\infty x_n^2/n$, zaś $f(x)=Q(x)-\|x\|_2^3$. Operator diagonalny $A:H\to H$, $(Ax)_n=x_n/n$, jest samosprzężony i ograniczony, ponieważ $\|Ax\|_2^2\le\|x\|_2^2$. Zatem $Q(x)=\langle Ax,x\rangle/2$ jest gładką funkcją kwadratową. Poniżej uzasadniamy, że $G(x)=\|x\|_2^3$ jest klasy $C^2$, a więc również $f$ ma tę regularność. Mamy $Df(0)=0$ i $D^2f(0)[h,h]=\sum_{n\ge1}h_n^2/n>0$ dla $h\ne0$, ale $0$ nie jest minimum lokalnym. Istotnie, dla $t_n=1/n$ mamy $t_ne_n\to0$ w normie i $f(t_ne_n)=1/(2n^3)-1/n^3=-1/(2n^3)<0=f(0)$. Funkcja $G(x)=\|x\|_2^3$ jest klasy $C^2$ również w zerze. Dla $x\ne0$ mamy
$DG(x)h=3\|x\|\langle x,h\rangle$ oraz
$D^2G(x)[h,k]=3\|x\|\langle h,k\rangle+3\langle x,h\rangle\langle x,k\rangle/\|x\|$.
Stąd $\|D^2G(x)\|\le6\|x\|$. Ponieważ $|G(h)-G(0)|/\|h\|=\|h\|^2\to0$, mamy $DG(0)=0$. Ponadto $\|DG(x)\|\le3\|x\|^2$, co dowodzi $D^2G(0)=0$ i ciągłości drugiej pochodnej w zerze.
\end{przyklad}

\subsection{Szereg Neumanna i odwracanie operatorów}

\begin{lemat}[Szereg Neumanna]
\label{lem:neumann15}
Niech $E$ będzie przestrzenią Banacha nad $\K\in\{\R,\C\}$, $I$ operatorem identycznościowym i $T\in\mathcal L(E)$ operatorem liniowym ograniczonym spełniającym $\|T\|<1$. Wtedy $I-T$ jest odwracalny oraz
$(I-T)^{-1}=\sum_{n=0}^\infty T^n$ w normie operatorowej, przy czym $\|(I-T)^{-1}\|\le(1-\|T\|)^{-1}$.
\end{lemat}
\begin{proof}
Szereg jest absolutnie zbieżny w przestrzeni Banacha $\mathcal L(E)$. Dla sum częściowych $S_N=\sum_{n=0}^N T^n$ zachodzi $(I-T)S_N=S_N(I-T)=I-T^{N+1}$; przechodząc do granicy, otrzymujemy obustronną odwrotność. Oszacowanie normy wynika z szeregu geometrycznego.
\end{proof}

\begin{stwierdzenie}[Otwartość grupy operatorów odwracalnych]
\label{prop:inverse-operator15}
Niech $E$ będzie rzeczywistą przestrzenią Banacha i $\mathrm{GL}(E)$ zbiorem odwracalnych operatorów w $\mathcal L(E)$. Wówczas $\mathrm{GL}(E)$ jest otwarty, a $\iota:\mathrm{GL}(E)\to\mathcal L(E)$, $\iota(A)=A^{-1}$, jest klasy $C^\infty$. Jej pochodna ma postać
$D\iota(A)H=-A^{-1}HA^{-1}$.
\end{stwierdzenie}
\begin{proof}
Dla $A\in\mathrm{GL}(E)$ i $H\in\mathcal L(E)$ spełniającego $\|A^{-1}H\|<1$ mamy $A+H=A(I+A^{-1}H)$, więc Lemat~\ref{lem:neumann15} daje odwracalność oraz lokalną jednostajną ograniczoność odwrotności. Z tożsamości
$(A+H)^{-1}-A^{-1}=-(A+H)^{-1}HA^{-1}$ wynika
\[
(A+H)^{-1}-A^{-1}+A^{-1}HA^{-1}
=A^{-1}H(A+H)^{-1}HA^{-1}.
\]
Dla $\|A^{-1}H\|\le1/2$ szereg Neumanna daje $\|(A+H)^{-1}\|\le2\|A^{-1}\|$. Norma prawej strony nie przekracza zatem $2\|A^{-1}\|^3\|H\|^2$, co dowodzi istnienia wskazanej pochodnej Frécheta. Jej ciągłość otrzymujemy z ciągłości odwracania i ciągłości kompozycji operatorów. Powtarzając różniczkowanie wzoru $D\iota(A)H=-A^{-1}HA^{-1}$ z regułą wieloliniową, otrzymujemy rekurencyjnie pochodne wszystkich rzędów i ich ciągłość.
\end{proof}

\subsection{Twierdzenie o funkcji odwrotnej}

\noindent\emph{Powiązanie tematyczne w Bibliotece matematycznej:} \atlasref{DG-13::thm:wspolrzedne-normalne}{DG-13: współrzędne normalne jako zastosowanie w skończonym wymiarze}.

\begin{twierdzenie}[Twierdzenie o funkcji odwrotnej dla przestrzeni Banacha]
\label{thm:inverse15}
Niech $E,F$ będą rzeczywistymi przestrzeniami Banacha, $U\subseteq E$ otwarty, $f\in C^1(U,F)$ i $x_0\in U$. Załóżmy, że $A=Df(x_0):E\to F$ jest izomorfizmem liniowym i topologicznym, tj. ma ograniczoną odwrotność $A^{-1}:F\to E$. Wówczas istnieją otoczenia otwarte $U_0\subseteq U$ punktu $x_0$ i $V_0\subseteq F$ punktu $y_0=f(x_0)$ takie, że $f|_{U_0}:U_0\to V_0$ jest bijekcją, a odwrotność $g:V_0\to U_0$ należy do $C^1(V_0,E)$ i spełnia
$Dg(y)=(Df(g(y)))^{-1}$ dla każdego $y\in V_0$.
\end{twierdzenie}
\begin{proof}
\emph{Krok 1: jednostajne przybliżenie liniowe.} Z ciągłości $Df$ w $x_0$ wybieramy $r>0$ tak małe, aby domknięta kula $\overline B(x_0,r)$ zawierała się w $U$ i
$\|I-A^{-1}Df(x)\|\le q<1$ dla $x\in\overline B(x_0,r)$; można przyjąć $q=1/2$. Dla $y\in F$ określamy $T_y(x)=x-A^{-1}(f(x)-y)$. Dla $x,z$ w tej kuli wzór całkowy na odcinku daje
\[
\|T_y(x)-T_y(z)\|\le q\|x-z\|.
\]
Ponadto $T_y(x_0)-x_0=A^{-1}(y-y_0)$. Dla $\|A^{-1}(y-y_0)\|\le(1-q)r$ operator $T_y$ przeprowadza $\overline B(x_0,r)$ w siebie, bo
$\|T_y(x)-x_0\|\le q\|x-x_0\|+\|A^{-1}(y-y_0)\|\le r$.

\emph{Krok 2: punkt stały.} Przypomnijmy zasadę kontrakcji: jeżeli $T$ jest odwzorowaniem $q$-lipschitzowskim pełnej przestrzeni metrycznej w siebie, $q<1$, to ciąg $x_{n+1}=T(x_n)$ jest Cauchy'ego, ponieważ $d(x_{n+1},x_n)\le q^n d(x_1,x_0)$. Dla $p>n$ nierówność trójkąta daje $d(x_p,x_n)\le q^n d(x_1,x_0)/(1-q)$, więc z zupełności ciąg ma granicę $x_*$. Ciągłość $T$ daje $T(x_*)=\lim_n T(x_n)=\lim_n x_{n+1}=x_*$. Dwa punkty stałe $u,v$ spełniają $d(u,v)\le qd(u,v)$, zatem są równe. Stosujemy tę zasadę do $T_y$ na domkniętej kuli w $E$. Otrzymujemy jedyny $x$ z $T_y(x)=x$, co jest równoważne $f(x)=y$. Jeżeli $\|A^{-1}(y-y_0)\|<(1-q)r$, to dla punktu stałego $x$ zachodzi $(1-q)\|x-x_0\|\le\|A^{-1}(y-y_0)\|<(1-q)r$, więc $x$ leży we wnętrzu kuli. Dodatkowo każdy $x\in B(x_0,r)$ ma na tym obszarze co najwyżej jeden punkt o tej samej wartości $f(x)$: jeżeli $f(x)=f(z)$, to $x,z$ są punktami stałymi tego samego odwzorowania $T_{f(x)}$, dla którego oszacowanie Lipschitza daje $\|x-z\|\le q\|x-z\|$, a więc $x=z$.

\emph{Krok 3: otwarte otoczenia i ciągłość odwrotności.} Ustalmy $0<r_0<r$ i połóżmy $U_0=B(x_0,r_0)$ oraz $V_0=f(U_0)$. Z oszacowania
$\|(x-z)-A^{-1}(f(x)-f(z))\|\le q\|x-z\|$ wynika
\[
\|x-z\|\le\frac{\|A^{-1}\|}{1-q}\,\|f(x)-f(z)\| \qquad(x,z\in B(x_0,r)).
\]
Zatem $f|_{U_0}$ jest różnowartościowe, a odwrotność $g:V_0\to U_0$ jest Lipschitzowska ze stałą $\|A^{-1}\|/(1-q)$. Zbiór $V_0$ jest otwarty: dla dowolnego $x\in U_0$ wybieramy kulę $\overline B(x,\delta)\subset U_0$; dla $\|A^{-1}(y-f(x))\|<(1-q)\delta$ i $z\in\overline B(x,\delta)$ zachodzi $\|T_y(z)-x\|\le q\|z-x\|+\|A^{-1}(y-f(x))\|<\delta$. Zasada kontrakcji daje rozwiązanie $f(z)=y$ w tej kuli. Warunek na $y$ określa otwarte otoczenie $f(x)$. W konsekwencji otoczenie $f(x)$ zawiera się w $V_0$.

\emph{Krok 4: różniczkowalność odwrotności.} Z wyboru $r$ mamy $\|I-A^{-1}Df(x)\|\le q<1$ dla $x\in U_0$. Lemat~\ref{lem:neumann15} daje odwracalność $A^{-1}Df(x)$ w $\mathcal L(E)$, a z $Df(x)=A(A^{-1}Df(x))$ wynika odwracalność $Df(x):E\to F$. Ustalmy $y=f(x)\in V_0$ i rozważmy $y+k=f(x+h)$, gdzie $h=g(y+k)-g(y)\to0$. Z różniczkowalności $f$ w $x$:
$k=Df(x)h+\rho(h)$ z $\|\rho(h)\|=o(\|h\|)$. Po pomnożeniu przez $(Df(x))^{-1}$ dostajemy
$h-(Df(x))^{-1}k=-(Df(x))^{-1}\rho(h)=o(\|h\|)=o(\|k\|)$,
gdyż otrzymana w kroku 3 nierówność zapewnia $\|h\|\le C\|k\|$. Stąd $Dg(y)=(Df(x))^{-1}$. Aby uzasadnić ciągłość odwracania także dla operatorów między różnymi przestrzeniami, oznaczmy przez $C=Df(x)$ i $B=Df(z)$ dla $z$ bliskich $x$. Mamy $B^{-1}=(I+C^{-1}(B-C))^{-1}C^{-1}$, a szereg Neumanna daje $B^{-1}\to C^{-1}$ przy $B\to C$ w normie operatorowej. Ciągłość $Df$ i $g$ daje więc ciągłość $Dg$. Tak więc $g\in C^1$.
\end{proof}

\begin{wniosek}[Wersja $C^m$]
\label{cor:inverse-ck15}
Niech $E,F$ będą rzeczywistymi przestrzeniami Banacha, $U\subseteq E$ otwarty, $x_0\in U$, a $f:U\to F$ odwzorowaniem klasy $C^m$, gdzie $m\ge2$, albo klasy $C^\infty$. Jeżeli $Df(x_0):E\to F$ jest izomorfizmem liniowym i topologicznym, to lokalna odwrotność $g$ z Twierdzenia~\ref{thm:inverse15} ma tę samą regularność co $f$.
\end{wniosek}
\begin{proof}
Ustalmy $A=Df(x_0)$. Wzór $Dg(y)=(A^{-1}Df(g(y)))^{-1}A^{-1}$ sprowadza odwracanie do $\mathrm{GL}(E)$, gdzie Stwierdzenie~\ref{prop:inverse-operator15} zapewnia gładkość. Reguła łańcuchowa wyższych rzędów w potrzebnej tu postaci wynika indukcyjnie z reguły pierwszego rzędu: pochodna złożenia jest iloczynem dwóch odwzorowań operatorowych, a ciągłe mnożenie dwuliniowe zachowuje klasę $C^j$ przez wielokrotne stosowanie reguły iloczynu. Możemy zatem prowadzić indukcję: jeśli $g\in C^j$ i $f\in C^{j+1}$, to $Df\circ g\in C^j$, zatem także $Dg\in C^j$, a więc $g\in C^{j+1}$. Zaczynamy od $g\in C^1$.
\end{proof}

\begin{przyklad}[Nieliniowe odwzorowanie na przestrzeni Hilberta]
\label{ex:inverse-hilbert15}
Niech $H$ będzie rzeczywistą przestrzenią Hilberta, $v\in H$ ustalony i $f:H\to H$ dane przez $f(x)=x+\|x\|^2v$. Wtedy $Df(x)h=h+2\langle x,h\rangle v$, a $Df(0)=I$. Twierdzenie~\ref{thm:inverse15} zapewnia lokalną odwrotność w otoczeniu zera. Jeżeli $v\ne0$, odwzorowanie nie jest nawet iniektywne na całym $H$: $f(0)=0$ i $f(-v/\|v\|^2)=0$, choć $-v/\|v\|^2\ne0$.
\end{przyklad}

\subsection{Twierdzenie o funkcji uwikłanej}

\begin{twierdzenie}[Twierdzenie o funkcji uwikłanej dla przestrzeni Banacha]
\label{thm:implicit15}
Niech $X,Y,Z$ będą rzeczywistymi przestrzeniami Banacha, $\mathcal U\subseteq X\times Y$ zbiorem otwartym, $F\in C^1(\mathcal U,Z)$, $(x_0,y_0)\in\mathcal U$ i $F(x_0,y_0)=0$. Załóżmy, że pochodna częściowa $D_yF(x_0,y_0):Y\to Z$ jest izomorfizmem liniowym i topologicznym. Wtedy istnieją otwarte otoczenia $U_0\subseteq X$ punktu $x_0$, $V_0\subseteq Y$ punktu $y_0$ takie, że $U_0\times V_0\subseteq\mathcal U$, oraz jedyna funkcja $\varphi:U_0\to V_0$ klasy $C^1$ taka, że dla $(x,y)\in U_0\times V_0$ zachodzi
$F(x,y)=0\ \Longleftrightarrow\ y=\varphi(x)$. Ponadto
\[
D\varphi(x)=-\bigl(D_yF(x,\varphi(x))\bigr)^{-1}D_xF(x,\varphi(x))\qquad(x\in U_0).
\]
Jeżeli $F\in C^m$ dla $m\ge2$ lub $C^\infty$, to również $\varphi$ ma odpowiednią regularność.
\end{twierdzenie}
\begin{proof}
Definiujemy $\Psi:\mathcal U\to X\times Z$ wzorem $\Psi(x,y)=(x,F(x,y))$. Jest to odwzorowanie klasy $C^1$ i
$D\Psi(x_0,y_0)(h,k)=(h,Bh+Ak)$, gdzie $A=D_yF(x_0,y_0)$ i $B=D_xF(x_0,y_0)$. Operator ten jest odwracalny: odwrotny przeprowadza $(h,z)$ w $(h,A^{-1}(z-Bh))$ i jest ograniczony. Twierdzenie~\ref{thm:inverse15} daje lokalny dyfeomorfizm $\Psi$ między otoczeniami $(x_0,y_0)$ oraz $(x_0,0)$. Oznaczmy dziedzinę i obraz lokalnego dyfeomorfizmu przez $\mathcal O$ i $\mathcal P$. Wybieramy otwarte otoczenia $U_1,V_0,W_0$ odpowiednio punktów $x_0,y_0,0$ takie, że $U_1\times V_0\subseteq\mathcal O$ i $U_1\times W_0\subseteq\mathcal P$; istnienie wynika z topologii produktu i otwartości $\mathcal O,\mathcal P$. Zapiszmy odwrotność $\Psi^{-1}(x,z)=(x,G(x,z))$; pierwsza współrzędna ma tę postać, ponieważ $\Psi$ zachowuje $x$. Ponieważ $G(x_0,0)=y_0$ i $G$ jest ciągła, wybieramy otwarte $U_0\subseteq U_1$ zawierające $x_0$ tak, by $G(x,0)\in V_0$ dla każdego $x\in U_0$. Połóżmy $\varphi(x)=G(x,0)$. Jest to funkcja klasy $C^1$, jako złożenie odwrotności z ciągłymi odwzorowaniami afinicznymi i rzutem. Wówczas $F(x,\varphi(x))=0$, a $\Psi$ jest iniektywna na $U_0\times V_0\subseteq\mathcal O$. Jeśli $F(x,y)=0$, to $\Psi(x,y)=(x,0)=\Psi(x,\varphi(x))$, zatem $y=\varphi(x)$. To daje równoważność oraz jednoznaczność funkcji $\varphi$, ponieważ dla każdego $x$ istnieje dokładnie jeden taki $y\in V_0$. Różniczkując tożsamość $F(x,\varphi(x))=0$ i używając reguły łańcuchowej, otrzymujemy
$D_xF(x,\varphi(x))+D_yF(x,\varphi(x))D\varphi(x)=0$.
Po ewentualnym zmniejszeniu $U_0$, z ciągłości $D_yF$ i $\varphi$ wynika $\|A^{-1}(D_yF(x,\varphi(x))-A)\|<1$; szereg Neumanna daje odwracalność $D_yF(x,\varphi(x)):Y\to Z$. Stąd wzór na $D\varphi$. Wersja wyższej regularności wynika z Wniosku~\ref{cor:inverse-ck15}.
\end{proof}

\begin{przyklad}[Rodzina równań operatorowych]
\label{ex:implicit-operator15}
Niech $E$ będzie rzeczywistą przestrzenią Banacha, $K\in\mathcal L(E)$ i $g\in E$. Na $\R\times E$ definiujemy $F(\lambda,u)=u-\lambda Ku-g$. W punkcie $(0,g)$ mamy $F(0,g)=0$ oraz $D_uF(0,g)=I$, zatem istnieje lokalnie jednoznaczne rozwiązanie $u(\lambda)$ klasy $C^\infty$ równania $(I-\lambda K)u(\lambda)=g$. Z wzoru na pochodną funkcji uwikłanej wynika
$u'(\lambda)=(I-\lambda K)^{-1}Ku(\lambda)$, w szczególności $u'(0)=Kg$. Różniczkując tożsamość $(I-\lambda K)u'(\lambda)=Ku(\lambda)$, otrzymujemy $(I-\lambda K)u''(\lambda)=2Ku'(\lambda)$, więc $u''(0)=2K^2g$. Dla $|\lambda|\,\|K\|<1$ rozwiązanie można również zapisać w postaci zbieżnego szeregu Neumanna $u(\lambda)=\sum_{n\ge0}\lambda^nK^ng$.
\end{przyklad}


\subsection{Różniczkowanie odwracania i eksponenty operatorowej}

\begin{twierdzenie}[Pochodne odwzorowania odwracania]
\label{thm:inverse-derivatives15}
Niech $X$ będzie przestrzenią Banacha nad $\K\in\{\R,\C\}$,
a $GL(X)$ zbiorem odwracalnych elementów algebry
$\mathcal L(X)$. Odwzorowanie $\operatorname{inv}(A)=A^{-1}$
jest klasy $C^\infty$ na $GL(X)$, przy czym w przypadku zespolonym różniczkujemy w sensie zespolonym Frécheta (kolejne pochodne utożsamiamy z operatorami wieloliniowymi nad $\K$). Dla każdego $A\in GL(X)$, $m\ge1$ i $H_1,\ldots,H_m\in\mathcal L(X)$ zachodzi
\[
 D^m\operatorname{inv}(A)[H_1,\ldots,H_m]
 =(-1)^m\sum_{\sigma\in S_m}
 A^{-1}H_{\sigma(1)}A^{-1}\cdots H_{\sigma(m)}A^{-1}.
\]
Tu $S_m$ oznacza zbiór permutacji $\{1,\ldots,m\}$.
\end{twierdzenie}
\begin{proof}
Jeżeli $\|A^{-1}H\|<1$, szereg Neumanna zapewnia $A+H\in GL(X)$, co dowodzi otwartości $GL(X)$ nad obydwoma ciałami, oraz daje
$(A+H)^{-1}=\sum_{k\ge0}(-A^{-1}H)^k A^{-1}$.
Suma wyrazów od $k=2$ ma normę co najwyżej
$\|A^{-1}\|\|A^{-1}H\|^2/(1-\|A^{-1}H\|)$.
Stąd ciągłość odwrotności i wzór
$D\operatorname{inv}(A)[H]=-A^{-1}HA^{-1}$.
Pochodna jest ciągła w normie operatorowej, dzięki
ciągłości mnożenia w algebrze. Indukcyjnie różniczkujemy
każdy czynnik $A^{-1}$ w formule rzędu $m$ regułą
iloczynu. Wstawienie $-A^{-1}H_{m+1}A^{-1}$ w miejsce
jednego z $m+1$ czynników $A^{-1}$ odpowiada wstawieniu
indeksu $m+1$ w jedną z $m+1$ pozycji permutacji.
Każda permutacja $m+1$ indeksów powstaje dokładnie raz.
Otrzymana suma jest ciągła jako odwzorowanie o wartościach
w przestrzeni ograniczonych operatorów wieloliniowych,
bo normę każdego składnika ogranicza iloczyn norm.
To dowodzi indukcyjnie wzoru i gładkości.
\end{proof}

\begin{twierdzenie}[Pochodne eksponenty operatorowej]
\label{thm:exp-derivatives15}
Niech $X$ będzie przestrzenią Banacha nad $\K\in\{\R,\C\}$.
Eksponenta $\exp A=\sum_{n\ge0}A^n/n!$ na $\mathcal L(X)$
jest klasy $C^\infty$ (w sensie zespolonym Frécheta, gdy $\K=\C$). Dla każdego $A\in\mathcal L(X)$, $m\ge1$ i $H_1,\ldots,H_m\in\mathcal L(X)$ zachodzi
\[
 D^m\exp(A)[H_1,\ldots,H_m]
 =\sum_{\sigma\in S_m}\int_{\Delta_m}
 e^{t_0 A} H_{\sigma(1)}e^{t_1 A}\cdots
 H_{\sigma(m)}e^{t_m A}\,\mathrm dt_1\cdots\mathrm dt_m,
\]
gdzie $S_m$ jest zbiorem permutacji $\{1,\ldots,m\}$, $\Delta_m=\{(t_1,\ldots,t_m)\in\R^m:t_j\ge0,\,
\sum_{j=1}^m t_j\le1\}$ i $t_0=1-\sum_{j=1}^m t_j$.
Całkowanie odbywa się względem $m$-wymiarowej miary Lebesgue'a ograniczonej do sympleksu. Można zdefiniować tę całkę przez iterowane całki Riemanna po zmiennych $t_1,\ldots,t_m$ z granicami wynikającymi z nierówności $\sum_jt_j\le1$; w dowodzie wykorzystujemy tę postać i całkowanie jednostajnie zbieżnego szeregu ciągłych funkcji. W szczególności dla $H\in\mathcal L(X)$
$D\exp(A)[H]=\int_0^1 e^{(1-t)A}He^{tA}\,\mathrm dt$.
\end{twierdzenie}
\begin{proof}
Pochodna rzędu $m$ wielomianu $A^n$ jest sumą wszystkich
słów z $n-m$ czynnikami $A$ oraz po jednym czynniku
$H_1,\ldots,H_m$ w dowolnej kolejności. Wynika to
indukcyjnie z reguły iloczynu: w każdym słowie zastępujemy
jeden z pozostałych czynników $A$ nowym kierunkiem.
Dla $n<m$ pochodna jest zerowa, a dla $n\ge m$ liczba słów wynosi $n(n-1)\cdots(n-m+1)=n!/(n-m)!$. Na kuli $\|A\|\le R$, gdzie $R>0$, i dla kierunków $\|H_j\|\le1$ norma każdego słowa nie przekracza $R^{n-m}$. Stąd norma pochodnej $A^n/n!$ jest nie większa niż $R^{n-m}/(n-m)!$.
Zatem dla każdego $m\ge0$ szereg pochodnych zbiega jednostajnie w normie operatorów wieloliniowych na każdej ograniczonej kuli, do ciągłej funkcji $P_m$; dla $m=0$ jest to szereg eksponenty. Jeżeli $P_{m,N}$ oznacza odpowiednią pochodną sumy częściowej, to wzór całkowy na przyrost (w przypadku zespolonym stosowany do krzywej z rzeczywistym parametrem) daje
\[
 P_{m,N}(A+H)-P_{m,N}(A)=\int_0^1 P_{m+1,N}(A+tH)[H,\cdot]\,\mathrm dt .
\]
Dla $m=0$ zapis $[H,\cdot]$ oznacza po prostu ewaluację na $H$. Jednostajna zbieżność na kuli zawierającej odcinek pozwala przejść do granicy. Po odjęciu $P_{m+1}(A)[H,\cdot]$ norma reszty nie przekracza $\|H\|\sup_{0\le t\le1}\|P_{m+1}(A+tH)-P_{m+1}(A)\|=o(\|H\|)$. Ponieważ $P_{m+1}(A)$ jest wieloliniowy nad $\K$, otrzymujemy pochodną Frécheta nad tym samym ciałem. Indukcja po $m$ uzasadnia wszystkie pochodne.
Grupując słowa według liczby $k_j$ czynników $A$
pomiędzy kierunkami, dostajemy
\[
 \sum_{\sigma\in S_m}\sum_{k_0,\ldots,k_m\ge0}
 \frac{A^{k_0}H_{\sigma(1)}A^{k_1}\cdots H_{\sigma(m)}A^{k_m}}
 {(m+k_0+\cdots+k_m)!}.
\]
Dla ustalonej permutacji rozwinięcie eksponent w całce na sympleksie jest jednostajnie bezwzględnie zbieżne w normie: supremum normy składnika o indeksach $(k_0,\ldots,k_m)$ nie przekracza $\prod_j\|H_j\|\prod_{i=0}^m\|A\|^{k_i}/k_i!$, bo $0\le t_i\le1$. Suma tych ograniczeń wynosi $e^{(m+1)\|A\|}\prod_j\|H_j\|$. Uzasadnia to całkowanie szeregu wyraz po wyrazie.
Dla całkowitych $k_j\ge0$ całka jednomianu
$t_0^{k_0}\cdots t_m^{k_m}$ po $\Delta_m$ wynosi
$\prod_{j=0}^m k_j!/(m+\sum_j k_j)!$.
Istotnie, dla $m=1$ wzór wynika z całkowania przez
części całki $\int_0^1 t^p(1-t)^q\,\mathrm dt$,
redukującego $p$ o jeden i zwiększającego $q$ o jeden,
z przypadkiem $p=0$ równym $1/(q+1)$. Dokładniej,
$I(p,q)=pI(p-1,q+1)/(q+1)$ dla $p\ge1$, skąd
$I(p,q)=p!q!/(p+q+1)!$.
W kroku indukcji ustalamy $t_m=u$ i podstawiamy
$t_j=(1-u)v_j$ dla $1\le j<m$. Dla $0\le u<1$
jakobian wynosi $(1-u)^{m-1}$, a $t_0=(1-u)(1-\sum_{j=1}^{m-1}v_j)$.
Przypadek $u=1$ nie wpływa na całkę. Oznaczmy
$P=m-1+\sum_{j=0}^{m-1}k_j$.
Z założenia indukcyjnego całka po $\Delta_{m-1}$,
przed uwzględnieniem czynników zależnych od $u$,
wynosi $\prod_{j=0}^{m-1}k_j!/P!$.
Mnożymy ją przez $\int_0^1u^{k_m}(1-u)^P\,\mathrm du
=k_m!P!/(k_m+P+1)!$ i otrzymujemy
$\prod_{j=0}^{m}k_j!/(m+\sum_{j=0}^m k_j)!$.
Całkowanie rozwinięcia wyraz po wyrazie daje dokładnie
powyższy szereg pochodnej.
\end{proof}

\begin{przyklad}[Nieprzemienność w pochodnej eksponenty]
\label{ex:noncommuting-exp15}
W algebrze macierzy $2\times2$ nad $\R$ lub $\C$ weźmy $A=\operatorname{diag}(0,1)$ i $H=E_{12}$ (jedynka w pozycji $(1,2)$, pozostałe elementy zerowe). Wtedy $e^{(1-t)A}He^{tA}=e^t H$. Zatem
$D\exp(A)[H]=(e-1)H$, podczas gdy $e^AH=H$ i
$He^A=eH$. Jeżeli $AH=HA$, całkowany operator
jest stale równy $e^AH$, więc w tym szczególnym
przypadku $D\exp(A)[H]=e^AH$.
\end{przyklad}

\begin{uwaga}[Dlaczego przestrzenie Banacha i dlaczego pochodna Frécheta?]
\label{rem:limits-inverse15}
Dowód twierdzenia o funkcji odwrotnej wykorzystuje zupełność domkniętej kuli (zasada kontrakcji) oraz jednostajną kontrolę operatora $Df(x)$ w normie. Sama odwracalność pochodnej Gâteaux w punkcie, bez stosownych założeń o ciągłości i pełnej różniczkowalności, nie daje takiego wniosku. W ogólnych przestrzeniach lokalnie wypukłych istnieją odrębne twierdzenia o funkcji odwrotnej, lecz nie wolno przenosić powyższego dowodu na dowolną topologiczną przestrzeń wektorową przez formalne zastąpienie normy seminormami.
\end{uwaga}
